% $Id: introduction.tex 1784 2012-04-27 23:29:31Z nicolas.cardozo $
% !TEX root = main.tex

\chapter{Introduction}

\label{cha:introduction}

Transfer learning enables a model to leverage knowledge gained from one task and apply it to another improving significantly efficiency and reducing the amount of time and data required for training. This powerful technique is widely used in various fields, including autonomous driving, where machine learning models must adapt to different driving scenarios. In autonomous vehicle development, levels of autonomy are classified on a scale from 1 to 5, where 1 represents a vehicle with no autonomy and 5 signifies full automation. Transfer learning plays a crucial role in advancing these levels by allowing pre-trained models to integrate knowledge from previous driving experiences, enabling vehicles to improve their ability to recognize and respond to complex road conditions by minimizing the additional training needed. This is particularly important in the context of autonomous driving, where the ability to generalize learned features to new tasks is essential for ensuring safety and efficiency on the road\cite{tl_survery_2016}.

Even for humans, the concept of “driving well” is inherently subjective, making it difficult to clearly define the appropriate behavioral outputs for an autonomous driving system. Similarly, selecting the optimal input features for learning is not straightforward. Driving behavior is only loosely defined, as different responses to similar situations can be equally valid. In recent years, reinforcement learning (RL) has emerged as a powerful framework for solving complex decision-making problems\cite{Sutton2018}, particularly in environments where explicit modeling of dynamics is challenging. RL allows an agent to learn optimal policies through interaction with an environment, receiving feedback in the form of rewards or penalties. This framework has been widely adopted in autonomous driving systems, where real-time decision-making, such as lane changing, is critical for ensuring safe and efficient driving\cite{dpr_for_real_driving_systems}.

Recent studies have demonstrated the effectiveness of transfer learning in mobile robotics, particularly in improving collision avoidance and navigation  \cite{collision_avoidance_tl}, This study presents a novel human-in-the-loop  as an expert that enhances transfer learning in mobile robots by integrating human expertise into reinforcement learning processes. By utilizing a pre-trained model that encapsulates human knowledge, the approach accelerates learning efficiency and improves generalization capabilities through collision avoidance experiments, the researchers demonstrated that incorporating model expert insights enables autonomous robots to achieve faster convergence and superior performance in complex environments.

This achievement underscores the role of transfer learning in bridging expert intuition with artificial intelligence, allowing mobile robots to adapt more effectively to dynamic real-world conditions. Despite its advantages, transfer learning faces critical limitations particularly in terms of both training data availability and time constraints. The adaptation of pre-trained models to new tasks often requires extensive fine-tuning, which can be computationally expensive and time-consuming. As a result, the growing interest in transfer learning stems not only from the need to accelerate AI development but also from the challenge of optimizing learning processes to reduce training time and data dependency while improving model performance across diverse applications\cite{tl_survey_2019}. 

In the context of autonomous vehicles, lane changing is a dynamic task that involves both short-term decision-making and long-term planning, requiring the vehicle to evaluate multiple factors, such as the surrounding traffic, road conditions, and safety. These driving system where multi task is mandatory, the classical supervised learning methods are not applicable, RL however, provides a suitable approach to this problem by enabling the vehicle to learn from experience and adapt to diverse driving scenarios\cite{dpr_for_real_driving_systems}.

Transfer learning (TL), a complementary technique to RL, aims to leverage knowledge gained from one domain to improve learning performance in a related but different domain. In self-driving mobile bases, transfer learning can be particularly useful for accelerating the learning process by transferring knowledge from simpler driving tasks (e.g., lane keeping or mobile base following) to more complex tasks like lane changing. This reduces the need for large amounts of training data and computational resources while improving generalization across different driving environments \cite{TL_survey}.

A practical example of transfer learning can be observed in human driving adaptation. A person who has learned to drive in London, where vehicles follow a left-side driving paradigm, does not need to start from zero when transitioning to driving in Colombia, where the right-side paradigm is used. Instead, they can somehow transfer their existing knowledge of vehicle control, traffic dynamics, and situational awareness while gradually adjusting to the new road rules and spatial orientation. This adaptation process is analogous to how transfer learning enables machine learning models to leverage previously acquired knowledge and modify it for new contexts, reducing the need for extensive retraining while enhancing efficiency in learning new driving conditions \cite{transfer_learning_driving_paradigm_change}.

In the current study, we explore the application of transfer learning in the context of lane changing for autonomous mobile bases. We investigate how knowledge gained from one driving paradigm can be effectively transferred to another, enabling the mobile base to adapt to new environments and driving conditions. By leveraging pre-trained models and expert demonstrations, we aim to enhance the learning efficiency and performance of the mobile base in navigating dynamic scenarios.

\subsection{Problem definition}
\label{sec:problem_definition}
Here comes the problem definition.

\subsection{Objectives}
\subsubsection{General objective}
To analyze and demonstrate the effectiveness of transfer learning in optimizing model adaptation to lane change paradigm in a controlled environment of a mobile robot in a roundabout, reducing training time requirements and the simulation steps.

\subsubsection{Specific objectives}
\begin{enumerate}
    \item To evaluate the impact of transfer learning on an autonomous mobile system by analyzing how expert models demonstrations enhance adaptability to new driving environments, measured in terms of transfer ratio, mastery, and pedagogy metrics.
    \item To compare the training efficiency and performance of an actor utilizing transferred knowledge versus an agent learning both tasks from scratch, analyzing improvements in learning speed, adaptation, and overall task execution quality.
    \item To evaluate the performance of an agent trained with expert demonstrations in its original environment, by quantifying its ability to replicate expert-level behaviors and maintain efficiency and accuracy using mastery metrics.
\end{enumerate}

\subsection{methodology}
% Aqui va una generalidad de la solución, se validó haciendo tal experimento y se evaluó asi. Concretamente lo que no descubrió, un parrafo max.

\subsection{Outline}
% Aqui va un parrafo de lo que se va a ver en cada capitulo, no es necesario que sea muy largo, pero si que sea claro y conciso. Un parrafo por capitulo.


\endinput

